\documentclass[11pt]{article}

\usepackage[margin=1in]{geometry}
\usepackage{amsmath,amssymb}
\usepackage[hidelinks]{hyperref}
\setlength{\parindent}{0pt}
\setlength{\parskip}{0.65em}

\begin{document}

\begin{center}
{\Large Literature-Already-Answered Status Note}\\[0.3em]
{\large Enriched Kannan mappings and completeness}\\[0.4em]
\textbf{Status:} literature already answered; not a new run result
\end{center}

\section*{Source question}

Berinde and P\u acurar introduce a self-map $T$ of a normed space $X$ as
\emph{enriched Kannan} when, for some $a\in[0,1/2)$ and $k\geq0$,
\[
 \|k(x-y)+Tx-Ty\|
 \leq a\bigl(\|x-Tx\|+\|y-Ty\|\bigr)
 \qquad(x,y\in X).
\]
On PDF page 5, item 3 of the remark following Theorem 2, they ask:

\begin{quote}
Does the enriched Kannan mapping principle still characterize the metric
completeness of the ambient space?
\end{quote}

This is the question in arXiv:1909.02379 \cite{BerindePacurar}; the published
version is identified as Remark 2.3 by the later paper. Since the displayed
definition uses vector addition and scalar multiplication, the precise
ambient category here is normed spaces.

\section*{Explicit later answer}

Yu, Li, and Ji reproduce the same question as Question 1 on printed page
21582 (PDF page 3) and state on printed page 21583 that they give an
affirmative answer \cite{YuLiJi}. Their Theorem 7 begins on that page, and its
proof runs through printed page 21585 (PDF page 6). The proof starts from a
nonconvergent Cauchy sequence in an incomplete normed space and constructs a
fixed-point-free enriched Kannan self-map. Together with the complete-space
fixed-point theorem of Berinde--P\u acurar, this establishes the intended
equivalence:
\[
 X\text{ is Banach}
 \quad\Longleftrightarrow\quad
 \text{every enriched Kannan self-map }T:X\to X\text{ has a fixed point}.
\]

The supporting authors explicitly know that they are answering the source
question: their abstract names the Berinde--P\u acurar question, and the main
text labels it Question 1 immediately before presenting Theorem 7.

\section*{Statement qualification}

The printed statement of Theorem 7 omits the universal fixed-point-property
quantifier and literally says that if one map $S$ has a fixed point, then
the normed space is complete. That literal statement is false: a constant
self-map is a Kannan, hence enriched Kannan, map with a fixed point on any
nonempty incomplete normed space. The surrounding affirmative-answer sentence
and the contradiction proof instead establish the universal characterization
displayed above.

The intended reverse implication also follows immediately from the source
definition: ordinary Kannan maps are precisely the $k=0$ subclass of the
enriched class. Hence the fixed-point property for all enriched Kannan maps
implies Subrahmanyam's fixed-point property for all ordinary Kannan maps, which
characterizes metric completeness \cite{Subrahmanyam}.

\section*{Classification and scope}

The match is explicit and exact, so this packet belongs in
\path{solutions/literature_already_answered/}. No part of the original
normed-space completeness question remains open. This note does not claim the
same formulation for arbitrary metric spaces without an added convex or
linear structure.

The status search used the exact phrases ``enriched Kannan mapping principle''
and ``characterizes completeness,'' plus the source arXiv id. It located the
2024 AIMS article, whose abstract and Question 1 explicitly cite the original
problem and whose Theorem 7 proof supplies the answer.

\begin{thebibliography}{9}

\bibitem{BerindePacurar}
V. Berinde and M. P\u acurar,
\emph{Fixed point theorems for Kannan type mappings with applications to split
feasibility and variational inequality problems},
arXiv:1909.02379, 2019; published as
\emph{Kannan's fixed point approximation for solving split feasibility and
variational inequality problems}, J. Comput. Appl. Math. 386 (2021), 113217,
\href{https://doi.org/10.1016/j.cam.2020.113217}{doi:10.1016/j.cam.2020.113217}.

\bibitem{YuLiJi}
Y. Yu, C. Li, and D. Ji,
\emph{Fixed point theorems for enriched Kannan-type mappings and application},
AIMS Mathematics 9(8) (2024), 21580--21595,
\href{https://doi.org/10.3934/math.20241048}{doi:10.3934/math.20241048}.

\bibitem{Subrahmanyam}
P. V. Subrahmanyam,
\emph{Completeness and fixed-points},
Monatsh. Math. 80 (1975), 325--330,
\href{https://doi.org/10.1007/BF01472580}{doi:10.1007/BF01472580}.

\end{thebibliography}

\end{document}
